
\section{From Capability to Verified Work}

Historical project duration is not a physical constant. When execution, verification, retries, and parallelism change, the workflow must be recomputed from its current graph $G=(V,E)$.

\subsection{Expected attempts and task duration}

For task $i$, let $p_i$ be autonomous success probability per attempt and $m_i$ the maximum attempts before escalation. The expected number of attempts is
\begin{equation}
N_i^{\mathrm{attempt}}=\frac{1-(1-p_i)^{m_i}}{p_i},
\end{equation}
and escalation probability is $(1-p_i)^{m_i}$. With AI execution time $\tau_i^{AI}$, verification time $\tau_i^V$, human fallback $\tau_i^H$, authority latency $\ell_i^A$, and external latency $\ell_i^E$,
\begin{equation}
\label{eq:task-duration}
\bar\tau_i=N_i^{\mathrm{attempt}}(\tau_i^{AI}+\tau_i^V)+(1-p_i)^{m_i}\tau_i^H+\ell_i^A+\ell_i^E.
\end{equation}
The separation matters: faster inference does not shorten a statutory waiting period, board approval, procurement cycle, laboratory incubation, or physical shipment.

\subsection{Critical path and effective parallelism}

Total expected work and critical path are
\begin{equation}
W=\sum_{i\in V}\bar\tau_i,\qquad
CP=\max_{\pi\in\mathcal P(G)}\sum_{i\in\pi}\bar\tau_i.
\end{equation}
With $n$ workers and parallel-efficiency exponent $\eta\in(0,1]$,
\begin{equation}
\label{eq:workflow-time}
\boxed{T_{\mathrm{workflow}}\approx\max\left\{CP,\frac{W}{n^\eta}\right\}+H_{\mathrm{coord}}(n).}
\end{equation}
This is an Amdahl-style constraint: parallel workers cannot erase a serial critical path \citep{amdahl1967,gustafson1988}.

\subsection{Reliability and verified cost}

If per-attempt verified reliability is $r_i$, then with $m_i$ independent attempts
\begin{equation}
r_i^*=1-(1-r_i)^{m_i}.
\end{equation}
For required steps,
\begin{equation}
\ln R_{\mathrm{workflow}}=\sum_i\ln r_i^*-\Omega_{\mathrm{corr}},
\end{equation}
where $\Omega_{\mathrm{corr}}\ge0$ penalizes common-mode failures. Expected cost is calculated from attempts, verification, and fallback; cost per verified successful workflow is
\begin{equation}
C_{\mathrm{verified}}=\frac{C_{\mathrm{workflow}}}{R_{\mathrm{workflow}}}.
\end{equation}
Raw token price is therefore not the decision-relevant cost when retries, verification, and failure severity dominate.

\subsection{Controlled workflow result}

The reference fixture contains parallel research tasks, a serial synthesis/review/release chain, retries, verification, and coordination overhead. Its total expected work is approximately \SI{14.88}{\hour}, but the critical path is \SI{13.26}{\hour}. Moving from one to two workers reduces duration; adding further workers eventually increases estimated duration because the critical path remains and coordination overhead rises.

\begin{figure}[H]
\centering
\includegraphics[width=0.84\textwidth]{workflow_parallelism.pdf}
\caption{Controlled workflow duration as parallel workers increase. The critical path prevents unlimited compression; coordination eventually offsets additional parallel capacity.}
\label{fig:workflow-parallelism}
\end{figure}

\begin{table}[H]
\centering
\caption{Selected controlled workflow outcomes.}
\begin{tabular}{rrrr}
\toprule
Workers & Duration (h) & Critical path (h) & Coord. overhead (h) \\
\midrule
1 & 14.88 & 13.26 & 0.00 \\
2 & 13.44 & 13.26 & 0.18 \\
4 & 13.62 & 13.26 & 0.36 \\
8 & 13.80 & 13.26 & 0.54 \\
32 & 14.16 & 13.26 & 0.90 \\
64 & 14.34 & 13.26 & 1.08 \\
\bottomrule
\end{tabular}

\label{tab:workflow}
\end{table}

The separate authority-latency example in the package demonstrates an even stronger case: a 24-hour authorization task dominates a workflow despite rapid machine execution. These examples are mechanical demonstrations, not estimates of a particular organization.
